\documentclass[11pt]{article}
\usepackage[margin=1in]{geometry}
\usepackage{amsmath,amssymb,amsthm,hyperref,enumitem}
\hypersetup{colorlinks=true,linkcolor=blue,urlcolor=blue}
\newtheorem{theorem}{Theorem}
\newcommand{\R}{\mathbb R}
\title{The Kadar--Yu root-distribution conjecture}
\author{Principia Math In Progress}
\date{22 September 2026}
\begin{document}
\maketitle
\noindent\textbf{MathDB 372352. A proof of Morris--Martin Conjecture 4.5.1.}

\section*{Task 1: all heights and all partitions}
Fix an integer \(l\ge-1\) and a partition \(\lambda\vdash l+2\).
Let \(P_n^{(\lambda)}\) be the reduced monic polynomial sequence
of [1, Theorem 3.19], defined by the one-cup Gram determinants
and extended backwards by its Chebyshev recurrence. Put
\[
 N=l+4,\qquad p_k(x)=P_{N+k}^{(\lambda)}(x),\qquad
 p_{k+1}(x)=xp_k(x)-p_{k-1}(x).                         \tag{1}
\]
Here ``reduced'' means that consecutive polynomials are coprime.
The minus sign in (1) is the determinant recurrence
[1, (3.2.7)]; it also appears in [2, slide 17].
The isolated plus sign in [1, (3.1.1)] is inconsistent with
these determinants and with its own classical examples.

\begin{theorem}
For every such \(l,\lambda\), each \(p_k\), \(k\ge-1\), has
only simple real roots. The numbers of roots in
\(( -\infty,-2)\) and \((2,\infty)\) are separately
nondecreasing for \(k\ge0\), bounded, and eventually constant.
Moreover, the union of the roots of \(p_k\), \(k\ge0\),
that lie in \([-2,2]\) is dense in that interval.
\end{theorem}

We prove this using the established Gram-matrix structure,
without assuming the explicit partition-family formulae
conjectured elsewhere in [1].

\section*{A positive resolvent from one tail}
Choose the real orthonormal Specht-module basis used in
[1, Section 3.2]. At rank \(N\), the Gram matrix has the form
\[
                         G(x)=xI-A,\qquad A=A^T\in M_s(\R).
                                                               \tag{2}
\]
This is [1, (3.2.3)], with the constant matrix negated.
Choose one terminal basis vector, corresponding to
\(u_{l+3,l+4}b_j\), and let \(e\in\R^s\) be its unit
coordinate vector. Extending just this tail by one vertex gives
the partial Gram matrix
\[
                  \begin{pmatrix}G(x)&e\\e^T&x\end{pmatrix}.
                                                               \tag{3}
\]
The other tails remain at rank \(N\). The proved one-tail
factorization [1, Lemma 3.24 and (3.2.16)--(3.2.17)] says
that the ratio of the determinant of (3) to \(\det G(x)\)
equals \(p_1(x)/p_0(x)\). Its Schur complement therefore yields
\[
 \frac{p_1(x)}{p_0(x)}=x-e^T(xI-A)^{-1}e,
 \qquad
 \frac{p_{-1}(x)}{p_0(x)}=e^T(xI-A)^{-1}e.             \tag{4}
\]
These are identities of rational functions, so no invertibility
assumption at particular real values is imposed.

By the spectral theorem, the last rational function is
\[
 m(x)=\sum_{j=1}^{q}\frac{w_j}{x-\theta_j},\qquad
 \theta_1<\cdots<\theta_q,\quad w_j>0,\quad
 \sum_{j=1}^{q}w_j=1.                                  \tag{5}
\]
Only eigenvalues with a nonzero projection of \(e\) are included.
Thus (5) has exactly \(q\) distinct poles. Since \(p_0,p_1\)
are coprime, so are \(p_0,p_{-1}=xp_0-p_1\).
Monicity and (4)--(5) identify
\[
 p_0(x)=\prod_{j=1}^{q}(x-\theta_j),\qquad
 \deg p_{-1}=q-1,
                                                               \tag{6}
\]
with \(p_{-1}\) monic. In particular \(q\ge1\).
On each interval \((\theta_j,\theta_{j+1})\),
\[
 m'(x)=-\sum_{i=1}^{q}\frac{w_i}{(x-\theta_i)^2}<0.
\]
The endpoint limits are \(+\infty\) and \(-\infty\).
Consequently \(p_{-1}\) has one simple root in each of those
\(q-1\) intervals and no other roots. When \(q=1\),
\(p_{-1}=1\), which is square free. This proves the
otherwise exceptional index \(k=-1\), as well as \(k=0\).

\section*{A finite matrix with an arbitrary path attached}
Let
\[
 \mathcal K=\operatorname{span}\{e,Ae,A^2e,\ldots\},\qquad
 H=A|_{\mathcal K}.
\]
This is a reducing subspace for \(A\). Its dimension is \(q\),
and \(e\) is cyclic for \(H\). Indeed, its spectral projections
are the \(q\) nonzero mutually orthogonal vectors in (5);
polynomial interpolation in the distinct \(\theta_j\) produces
each projection from \(e\).

Define \(H_k\) by attaching a path of \(k\) new vertices to
\(e\) in \(H\). Every new diagonal entry is zero and every
path edge, including the attachment, has weight one.
Thus \(H_0=H\) and \(H_1=\left(\begin{smallmatrix}H&e\\e^T&0
\end{smallmatrix}\right)\).
Equation (4), followed by expansion at each new terminal vertex,
gives
\[
                         \det(xI-H_k)=p_k(x),\qquad k\ge0.
                                                               \tag{7}
\]
For the initial step, \(\det(xI-H)=p_0\) by (6), and
\(\det(xI-H_1)=p_0(x)(x-m(x))=p_1(x)\).
All subsequent determinants obey (1).

The terminal vector of the attached path is cyclic for \(H_k\).
Successive applications of \(H_k\), subtracting already recovered
coordinates, recover the whole path and then \(e\). Continuing
recovers every \(H^je\), hence the entire core \(\mathcal K\).
For \(k=0\), cyclicity was already established.
A real symmetric matrix with a cyclic vector has simple spectrum:
an eigenspace of dimension at least two would contain a nonzero
vector perpendicular to the cyclic vector, and therefore to all
its matrix iterates. Equation (7) now proves squarefreeness
and reality for every \(k\ge0\).

\section*{Exterior roots and density}
Let \(L_k\) and \(R_k\) count the eigenvalues of \(H_k\)
strictly below \(-2\) and strictly above \(2\), respectively.
The matrix \(H_k\) is a principal submatrix of \(H_{k+1}\).
Cauchy interlacing consequently gives
\[
 L_{k+1}\ge L_k,\qquad R_{k+1}\ge R_k.                 \tag{8}
\]
Deleting the \(q\) core coordinates from \(H_k\) leaves the
ordinary \(k\)-vertex path matrix \(T_k\). Its eigenvalues are
\[
                  2\cos\frac{j\pi}{k+1},\qquad 1\le j\le k,
                                                               \tag{9}
\]
as follows by substituting the sine vectors into its three-term
difference equation. They all lie in \((-2,2)\).
Interlacing for this deletion implies
\(L_k\le q\) and \(R_k\le q\).
Bounded nondecreasing integer sequences are eventually constant.
This also accounts for possible roots exactly at the endpoints:
the counts throughout are on the two open exterior intervals.

For completeness, interlacing gives a quantitative density result.
Let \(F_B(t)\) count the eigenvalues of \(B\) that are at most \(t\).
Deleting \(q\) coordinates yields
\[
                   F_{T_k}(t)\le F_{H_k}(t)\le F_{T_k}(t)+q.
                                                               \tag{10}
\]
Define the arcsine distribution function
\[
 F(t)=\begin{cases}
 0,&t\le-2,\\
 1-\pi^{-1}\arccos(t/2),&-2<t<2,\\
 1,&t\ge2.
 \end{cases}
\]
The explicit grid (9) shows
\(\lvert F_{T_k}(t)-kF(t)\rvert\le1\). Combining with
(10), and using \(\deg p_k=q+k\), gives
\[
 \sup_{t\in\R}\left|
 \frac{F_{H_k}(t)}{q+k}-F(t)\right|
 \le\frac{q+1}{q+k},\qquad k\ge1.                    \tag{11}
\]
Every nonempty open interval inside \((-2,2)\) has positive
arcsine measure. By (11), it contains roots of \(p_k\) for
all sufficiently large \(k\). Intervals arbitrarily near either
endpoint give density in the closed interval \([-2,2]\).
Together with (6)--(8), this proves all clauses of the conjecture.

\section*{Tasks 2--5: proposed mathematical follow-ups}
\begin{enumerate}[label=\textbf{Task \arabic*:},start=2,leftmargin=*]
\item Express the spectral points and weights in (5) directly
from the Young diagram of \(\lambda\), without constructing the
full one-cup Gram matrix.
\item Determine the eventual values of \(L_k,R_k\) and an
explicit stabilization index in terms of \(l,\lambda\).
\item Determine the limits and convergence rates of the exterior
roots, and characterize bounded generalized eigenvectors of the
infinite path extension at \(-2\) and \(2\).
\item Derive a phase-shift formula for individual interior roots
and sharp spacing estimates, refining the distribution bound (11).
\end{enumerate}

\begin{thebibliography}{9}
\bibitem{MM}
B.~Morris and P.~P.~Martin,
\emph{On semisimplicity criteria and non-semisimple representation
theory for the Kadar--Yu algebras},
\href{https://arxiv.org/abs/2512.24535}{arXiv:2512.24535v1}
(31 December 2025), 66 pages.
The Gram structure and one-tail factorization in Section 3
are established input; the target is Conjecture 4.5.1 on p.~35.
\bibitem{talk}
B.~Morris, \emph{Interpolating between Temperley--Lieb and Brauer
representation theory: generalised roots of unity},
\href{https://drive.google.com/file/d/1awcjFlsm4D7ApgIJ7WjTw_CcWoHwmTMp/view}
{ARTIN slides}, Newcastle University, 11 September 2026.
\end{thebibliography}
\end{document}
